\section{Верификация численного алгоритма}
\label{sec:verification}

\subsection{Метод изготовленных решений}

Для проверки численного алгоритма применяется метод изготовленных решений~\cite{salari2000}, в котором заранее задаётся температурное поле $T_*$, а источник, начальные и граничные условия выбираются так, чтобы это поле удовлетворяло уравнению теплопереноса~\eqref{eq:heat}. Подстановка $T_*$ в исходное уравнение определяет объёмную плотность мощности проверочного источника
\begin{equation}
 A_* = (\rho c)\left(\frac{\partial T_*}{\partial t}
       + \boldsymbol v \cdot \nabla T_*\right) - \nabla \cdot (\lambda\nabla T_*).
 \label{eq:mms-source}
\end{equation}
Решение полученной задачи тем же численным алгоритмом и его сопоставление с $T_*$ позволяют проверить аппроксимацию тепловых потоков, интегрирование по времени и учёт движения материала. Погрешность оценивается в центрах ячеек в относительной норме
\begin{equation}
 E_T = 100\left(\frac{\sum_P |P|\bigl(T_P - T_*(x_P,z_P,t)\bigr)^2}
 {\sum_P |P|\bigl(T_*(x_P,z_P,t) - 273{,}15\bigr)^2}\right)^{1/2}\quad[\%],
 \label{eq:mms-error}
\end{equation}
где веса $|P|$ учитывают площади ячеек, а нормирование по температуре в градусах Цельсия согласуется с оценками в прикладных расчётах.

\subsubsection{Стационарная задача в одном слое}

Для оценки пространственной аппроксимации рассматривается слой шириной $W=4000$~м и мощностью $h_0=2000$~м с постоянной либо линейно меняющейся теплопроводностью. В безразмерных координатах $\xi=x/W$ и $\eta=z/h(t)$ температура и коэффициенты задаются соотношениями
\begin{gather}
 T_* = 300 + 20\psi(t)\cos(\pi\xi)\sin(\pi\eta/2),\qquad
 \rho c = 3 \cdot 10^6,\label{eq:mms-layer}\\
 \lambda_* = 2\quad\text{либо}\quad 2(1 + 0{,}2\xi + 0{,}3\eta).
\end{gather}
Температура выражается в кельвинах, теплопроводность~--- в Вт/(м$\cdot$К), объёмная теплоёмкость~--- в Дж/(м$^3\cdot$К). В стационарной задаче принимаются $\psi=1$ и $h=h_0$, на кровле задаётся температура 300~К, а на остальных границах~--- нулевой поток. Выбранное тригонометрическое поле позволяет оценить погрешность аппроксимации по её уменьшению при сгущении сетки.

На рис.~\ref{fig:spatial} ошибка~\eqref{eq:mms-error} сопоставлена с квадратичной зависимостью от шага сетки. По отношению ошибок на двух последних сетках для MPFA получены порядки 2,14 при постоянной теплопроводности и 1,96 при переменной, а для NMPFA~--- 1,89 в обоих случаях, что подтверждает сходимость, близкую ко второму порядку.

\reportfigure{07_spatial_mms.png} {Сходимость пространственной аппроксимации на гладком стационарном решении, заданном методом изготовленных решений. (а) постоянная теплопроводность; (б) теплопроводность, линейно меняющаяся по координатам. Показаны относительные ошибки температуры для MPFA и NMPFA, чёрная линия соответствует второму порядку.} {fig:spatial}

\subsubsection{Контакт материалов с разной теплопроводностью}

Для проверки межслойных потоков рассматриваются два материала, разделённые прямой $z=z_c+bx$. Координата $s=z-z_c-bx$ обращается в нуль на контакте, а температура в материале $\ell\in\{1,2\}$ задаётся синусоидальной функцией с амплитудой, обратно пропорциональной его теплопроводности,
\begin{equation}
 T_* = 300 + B\frac{\lambda_1}{\lambda_\ell}\sin(\kappa s).
 \label{eq:mms-contact}
\end{equation}
Такой выбор обеспечивает непрерывность температуры и нормального теплового потока при $s=0$, а соотношение~\eqref{eq:mms-source} в стационарном случае определяет проверочный источник
\begin{equation}
 A_* = B\lambda_1\kappa^2(1 + b^2)\sin(\kappa s).
\end{equation}

В области размером $4\times2{,}6$~м при $z_c=0{,}7$~м сопоставляются горизонтальный контакт с $b=0$ и наклонный с $b=0{,}3$. Приняты $B=20$~К, $\kappa=\pi/(2{,}6\,\text{м})$, $\lambda_1=1$~Вт/(м$\cdot$К) и отношения теплопроводностей $\lambda_2/\lambda_1=100$ и 1000, а граничные условия определяются из функции~\eqref{eq:mms-contact}. При горизонтальном контакте обе схемы дают сходимость, близкую ко второму порядку, при наклонном~--- к первому, что показано сравнением с соответствующими степенными зависимостями на рис.~\ref{fig:contacts}.

\reportfigure{08_contact_mms.png} {Сходимость на стационарном синусоидальном решении с непрерывными температурой и тепловым потоком на контакте двух материалов. (а) горизонтальный контакт; (б) наклонный контакт. Сопоставлены относительные ошибки температуры для MPFA и NMPFA при отношениях теплопроводностей 100:1 и 1000:1. Чёрные линии соответствуют первому и второму порядкам.} {fig:contacts}

Чтобы установить, какие потоки определяют поведение ошибки на наклонном контакте, дополнительно рассмотрено кусочно-линейное поле с непрерывным тепловым потоком. При одинаковой теплопроводности материалов оно воспроизводится с точностью округления, а при различающейся теплопроводности сохраняется первый порядок сходимости. Подстановка точных температур в формулы потоков локализует расхождения на вертикальных подгранях между разными материалами, тогда как внутри одного материала потоки воспроизводятся с точностью округления.

По результатам этих численных экспериментов поведение ошибки на наклонном контакте определяется его ступенчатым представлением и аппроксимацией межслойных потоков. При фиксированном $\Delta x=0{,}03125$~м и отношении теплопроводностей 100:1 уменьшение $\Delta z/\Delta x$ от 1 до 0,25 снижает ошибку синусоидального решения примерно в 180~раз для MPFA и в 33~раза для NMPFA (табл.~\ref{tab:contact-diagnostics}), поэтому более мелкий вертикальный шаг повышает точность в рассматриваемой геометрии.

\begin{table}[htbp]
\centering\small
\caption{Влияние вертикального шага на ошибку температуры при наклонном контакте.}\label{tab:contact-diagnostics}
\begin{tabular}{lrr}
\toprule
Схема & $\Delta z/\Delta x$ & $E_T$, \% \\
\midrule
MPFA & 0,25 & 0,00117 \\
MPFA & 1,00 & 0,20759 \\
NMPFA & 0,25 & 0,00612 \\
NMPFA & 1,00 & 0,20278 \\
\bottomrule
\end{tabular}
\end{table}


\subsubsection{Нестационарная задача на неподвижной и подвижной сетках}

Для проверки интегрирования по времени и движения сетки в решении~\eqref{eq:mms-layer} принимаются переменная теплопроводность и временной множитель $\psi(t)=\exp(-t/t_{\mathrm{final}})$ при $t_{\mathrm{final}}=10^{12}$~с. Сопоставляются неподвижный слой с $h=h_0$ и слой с возрастающей мощностью $h=h_0(1+t/(4t_{\mathrm{final}}))$, причём начальная температура задаётся как $T_*(x,z,0)$, а скорость материала определяется движением границ.

В расчётах последовательно используются $n_x=16,32,64,128$~колонок и 8, 32, 128, 512~временных шагов. Число ячеек по вертикали в соседних колонках чередуется между $n_x/2$ и $n_x/2+1$, чтобы несовпадение уровней ячеек сохранялось при сгущении. При втором порядке пространственной и первом порядке временной аппроксимации ожидается поведение ошибки
\begin{equation}
 E_T = \mathcal O(\Delta x^2 + \Delta z^2 + \Delta t).
 \label{eq:mms-order}
\end{equation}
В соответствии с оценкой~\eqref{eq:mms-order} пространственные шаги уменьшаются вдвое, а временной~--- вчетверо, чтобы оба вклада в ошибку снижались приблизительно в четыре раза. На двух последних уровнях сгущения получено уменьшение ошибки в 3,99--4,00~раза для обеих схем на неподвижном и движущемся слоях (рис.~\ref{fig:time}), что согласуется с ожидаемыми порядками аппроксимации.

\reportfigure{09_joint_mms.png} {Уменьшение ошибки нестационарного изготовленного решения при сгущении сетки и временного шага. Сопоставлены MPFA и NMPFA для неподвижного и движущегося слоёв с переменной теплопроводностью. При уменьшении пространственных шагов вдвое временной шаг уменьшается вчетверо, чёрная линия соответствует квадратичной зависимости ошибки от пространственного шага.} {fig:time}

\subsection{Влияние минимальной высоты ячеек}

Минимальная высота ячеек $\varepsilon$ вводит дополнительное тепловое сопротивление в ещё не отложившемся слое, поэтому её влияние оценивается на одномерной задаче с постоянным базальным потоком. Верхняя часть неподвижного слоя заменяется прослойкой толщиной $\delta h=N\varepsilon$ с теплопроводностью $\lambda_{\mathrm{new}}$ вместо исходной $\lambda_{\mathrm{old}}$, вследствие чего стационарная температура под прослойкой изменяется на
\begin{equation}
 \Delta T = q\,N\varepsilon
 \left(\frac{1}{\lambda_{\mathrm{new}}}
             - \frac{1}{\lambda_{\mathrm{old}}}\right).
\end{equation}
В проверке приняты $N=40$, $q=75$~мВт/м$^2$, $\lambda_{\mathrm{new}}=1$~Вт/(м$\cdot$К) и $\lambda_{\mathrm{old}}=2$~Вт/(м$\cdot$К), а отражательная способность вычисляется за 50~млн~лет при полученной постоянной температуре.

При используемой высоте $\varepsilon=0{,}05$~м суммарная мощность прослойки составляет 2~м, а максимальные относительные различия с решением без неё равны 0,546\% для температуры и 0,140\% для отражательной способности при нормировании температурных различий по градусам Цельсия (табл.~\ref{tab:minimum}). Таким образом, в этой проверке ненулевая высота сжатых ячеек обеспечивает представление нового слоя при изменении расчётных полей менее чем на 1\%.

\begin{table}[htbp]
\centering\small
\caption{Влияние минимальной высоты ячейки; максимальные относительные отличия.}\label{tab:minimum}
\begin{tabular}{lrrr}
\toprule
$\varepsilon$, м & $N\varepsilon$, м & T, \% & Ro, \% \\
\midrule
0,001 & 0,04 & 0,0109 & 0,0028 \\
0,01 & 0,40 & 0,1091 & 0,0280 \\
0,025 & 1,00 & 0,2727 & 0,0699 \\
0,05 & 2,00 & 0,5455 & 0,1399 \\
0,1 & 4,00 & 1,0909 & 0,2800 \\
0,5 & 20,00 & 5,4545 & 1,4080 \\
\bottomrule
\end{tabular}
\end{table}


\subsection{Обратная задача с синтетическими наблюдениями}

Для проверки восстановления базального потока используются геометрия Модели~3 и наблюдения, полученные из прямого расчёта с известной историей $q(t)$. Заданный поток состоит из трёх линейных участков с узлами $(\tau,q)=(200,90)$, $(150,80)$, $(55,77)$ и $(0,60)$, где возраст выражен в млн~лет, а поток~--- в мВт/м$^2$. Из современных полей выбираются значения в тех же точках пяти скважин, что и в исходном наборе данных, и полученные 500~температур и 50~значений отражательной способности используются как синтетические наблюдения без случайных возмущений.

При подборе определяются все 17~параметров на возрастах конфигураций, а прямой расчёт и обратная задача решаются с одинаковыми шагами $\Delta x=400$~м, $\Delta z=100$~м и $\Delta t=1$~млн~лет. Невязки нормируются на 2\% от наблюдений, как в основном расчёте, но для данных без случайных возмущений задаётся более строгий критерий остановки $|r_j|\leq10^{-3}$. Начальная, заданная и восстановленная истории потока сопоставлены на рис.~\ref{fig:synthetic}.

В результате минимизации получено $J=0{,}00366$, а максимальные относительные невязки составили 0,0029\% для температуры и 0,0391\% для отражательной способности (табл.~\ref{tab:synthetic}). При этом относительное различие заданного и восстановленного потоков в ранней части истории достигает 11,01\%, а различие их интегралов по времени составляет 3,11\%.

\reportfigure{10_synthetic_inverse.png} {Восстановление базального потока по синтетическим наблюдениям на геометрии Модели~3. Чёрная линия обозначает заданную историю из трёх линейных участков, синяя~--- начальное приближение, красная~--- восстановленную историю. Точками отмечены все 17~узлов, соответствующих возрастам конфигураций.} {fig:synthetic}
\begin{table}[htbp]
\centering\small
\caption{Согласование синтетических наблюдений.}\label{tab:synthetic}
\begin{tabular}{lrrrr}
\toprule
Поле & Наблюдений & В пределах 2\% & $\varepsilon_{\mathrm{rms}}$, \% & Максимум, \% \\
\midrule
T & 500 & 500 & 0,00136 & 0,00286 \\
Ro & 50 & 50 & 0,01655 & 0,03910 \\
\bottomrule
\end{tabular}
\end{table}


Для оценки полей за пределами точек наблюдений сопоставлены прямые решения с заданным и восстановленным потоками во всех 17~конфигурациях. Поля определяются на исходной геометрии линейной интерполяцией между кровлей и подошвой слоёв и соседними горизонтальными узлами, а их различия оцениваются в относительной норме $L^2$ по площади разреза
\begin{equation}
 E_u(t) = 100\frac{\|u^{\mathrm{calc}} - u^{\mathrm{ref}}\|_{L^2(\Omega(t))}}
                    {\|u^{\mathrm{ref}}\|_{L^2(\Omega(t))}}\quad[\%],
 \qquad u \in \{T_C,R_o\},
 \label{eq:field-error}
\end{equation}
где индексы $\mathrm{ref}$ и $\mathrm{calc}$ обозначают решения с заданным и восстановленным потоками соответственно.

В современный момент нормы различий температуры и отражательной способности равны 0,00173 и 0,01870\%, а для более ранних конфигураций в целом возрастают (рис.~\ref{fig:synthetic-field-errors}). Максимальные значения составляют 5,63\% для температуры при возрасте 187,5~млн~лет и 1,22\% для отражательной способности при возрасте 55~млн~лет, пространственные распределения относительных разностей на этих возрастах и в современности показаны на рис.~\ref{fig:synthetic-temperature} и~\ref{fig:synthetic-ro}.

\reportfigure{11_synthetic_field_errors.png} {Различия полей температуры и отражательной способности во всех конфигурациях синтетической задачи в зависимости от геологического возраста. Сопоставлены прямые решения при заданном и восстановленном потоках в относительной норме $L^2$ по площади разреза. Для каждого поля нормирование выполнено по решению с заданным потоком.} {fig:synthetic-field-errors}
\reportfigure{12_synthetic_temperature.png} {Пространственное распределение относительной разности температур в синтетической задаче. (а) возраст 187,5~млн~лет; (б) возраст 55~млн~лет; (в) современное состояние. Показано отклонение решения с восстановленным потоком от решения с заданным потоком, нормированное по последнему в градусах Цельсия. Для всех панелей используются общие шкалы глубины и относительной разности.} {fig:synthetic-temperature}
\reportfigure{13_synthetic_reflectance.png} {Пространственное распределение относительной разности отражательной способности витринита в синтетической задаче. (а) возраст 187,5~млн~лет; (б) возраст 55~млн~лет; (в) современное состояние. Показано отклонение решения с восстановленным потоком от решения с заданным потоком, нормированное по последнему. Для всех панелей используются общие шкалы глубины и относительной разности.} {fig:synthetic-ro}

Таким образом, изготовленные решения подтверждают ожидаемое уменьшение ошибки на гладких полях и позволяют определить точность потоков на контактах материалов, а синтетические наблюдения воспроизводятся с относительными невязками менее 0,04\%. Восстановленный поток даёт близкие современные поля во всей области, при этом наибольшие различия с заданной температурной историей приходятся на ранние конфигурации.
